% !TITLE 唐
% !SUMMARY 中国诗歌的黄金时代
% !ORDER 0

\begin{intro}
如果中国诗歌史只能选择一个时代来代表，那无疑是唐代。

唐朝近三百年，政治开明、经济繁荣、文化包容，为诗歌的勃发提供了最肥沃的土壤。从初唐四杰到陈子昂，从盛唐的李杜到王维、孟浩然，从中唐的韩孟、元白到刘禹锡、柳宗元，再到晚唐的李商隐、杜牧——几乎每一个时段都有光芒四射的名字。全唐诗收录诗人两千余家、诗作四万八千余首，这还只是流传下来的部分。

唐代诗歌的伟大，首先在于它的气象。"黄河之水天上来，奔流到海不复回"——这是盛唐的声音，自信、豪迈、充满生命的张力。即使是写边塞苦寒、写离别相思，盛唐诗人也总能在悲怆中保持一份昂扬。王维的"大漠孤烟直，长河落日圆"，在极简的线条里勾勒出天地的壮阔；李白的"仰天大笑出门去，我辈岂是蓬蒿人"，把个人抱负写得如此张扬而不令人反感——这需要何等的气度。

杜甫则代表了另一种伟大。他经历了安史之乱，从盛唐的繁华跌落到乱世的流离，却用诗歌记录下了一个时代的创伤。"三吏""三别"是诗史，他的律诗更是达到了艺术的极致。"无边落木萧萧下，不尽长江滚滚来"——对仗之工、意境之深、情感之沉郁，后世几乎无人能及。

唐诗还完成了中国诗歌形式的最终定型。五七言古体诗在唐代达到了巅峰，而律诗、绝句的格律也在此时成熟。形式与内容的完美统一，使得唐诗成为中国文学史上最具辨识度的艺术形态。读唐诗，就是在读中国人情感表达的最精炼、最优美的方式。
\end{intro}
